% ── CAPÍTULO 1: INTRODUCCIÓN Y VISIÓN GENERAL ───────────────
\section{Introducción y Visión General}
\label{sec:intro}

\subsection{¿Qué es K-QGMIP?}

\textbf{K-QGMIP} es una herramienta de software desarrollada en Python para analizar sistemas de variables binarias con dinámicas estocásticas. Su objetivo principal es encontrar la \textbf{partición mínima de información} (\textit{k-MIP}) de un sistema: la forma de dividir las variables en $k$ grupos de modo que se pierda la menor cantidad posible de información al tratar cada grupo de forma independiente.

\begin{nota}
Este software está diseñado para investigadores e ingenieros interesados en la \textbf{Teoría de la Información Integrada (IIT)}: un marco matemático que estudia cómo los sistemas integran información y que está relacionado con el estudio de la conciencia en sistemas físicos.
\end{nota}

\subsection{¿Para qué sirve?}

El software responde a la pregunta: \textit{¿cuál es la mejor manera de dividir las $n$ variables de un sistema en exactamente $k$ grupos, de forma que la pérdida de información sea mínima?}

\vspace{6pt}
\textbf{Aplicaciones típicas:}
\begin{itemize}[itemsep=4pt]
    \item Análisis de redes neuronales artificiales o biológicas.
    \item Estudio de dependencias entre variables en sistemas complejos.
    \item Comparación de múltiples estrategias de particionamiento ($k=2$ hasta $k=5$).
    \item Exploración de cómo aumentar $k$ reduce la pérdida de información.
\end{itemize}

\subsection{Estrategias Disponibles}

K-QGMIP implementa cuatro estrategias (dos bases y dos extensiones):

\begin{table}[H]
    \centering
    \caption{Estrategias de particionamiento disponibles.}
    \begin{tabular}{lll}
        \toprule
        \textbf{Estrategia} & \textbf{k soportado} & \textbf{Descripción breve} \\
        \midrule
        GeoMIP        & $k=2$ & Búsqueda geométrica en hipercubo \\
        KGeoMIP       & $k=2\ldots5$ & Extensión de GeoMIP a $k$ grupos \\
        QNodes        & $k=2$ & Heurística MAO con oráculo submodular \\
        KQNodes       & $k=2\ldots5$ & Extensión de QNodes a $k$ grupos \\
        \bottomrule
    \end{tabular}
\end{table}

\subsection{Resultado Principal: el Dashboard}

Al terminar el benchmark, el software genera un archivo \textbf{dashboard.html} que puede abrirse en cualquier navegador web. Contiene:

\begin{enumerate}[label=\textbf{\arabic*.}]
    \item \textbf{Tiempos de ejecución} por instancia y estrategia (gráfico de líneas interactivo).
    \item \textbf{Correlación de pérdidas} entre estrategias pares (diagramas de dispersión).
    \item \textbf{Distribución estadística} de tiempos (boxplots).
\end{enumerate}

\subsection{Convenciones de Este Manual}

\begin{itemize}[itemsep=4pt]
    \item \lstinline|texto así| indica comandos de terminal o nombres de archivos/carpetas.
    \item \textbf{Negrita} marca conceptos importantes la primera vez que aparecen.
    \item Las cajas de colores señalan información especial:
\end{itemize}

\begin{nota}
Información útil o contexto adicional.
\end{nota}
\begin{aviso}
Advertencia: algo podría salir mal si no se sigue la indicación.
\end{aviso}
\begin{peligro}
Acción que puede dañar datos o generar resultados incorrectos.
\end{peligro}
\begin{tip}
Sugerencia para trabajar de forma más eficiente.
\end{tip}
